% ═══════════════════════════════════════════════════════════════
% LERNBEGLEITER: Waagerechter Wurf (Physik Klasse 10E)
% 2 Seiten A4 – Begleitmaterial zum digitalen Lernpfad
% ═══════════════════════════════════════════════════════════════

\documentclass[10pt, a4paper]{article}

\usepackage[utf8]{inputenc}
\usepackage[T1]{fontenc}
\usepackage[ngerman]{babel}
\usepackage{helvet}
\renewcommand{\familydefault}{\sfdefault}

\usepackage[margin=1.2cm, top=1.5cm, bottom=1.2cm]{geometry}
\usepackage{enumitem}
\usepackage{tabularx}
\usepackage{array}
\usepackage{multicol}
\usepackage{amsmath}
\usepackage{amssymb}

\usepackage{xcolor}
\usepackage{tcolorbox}
\usepackage{fancyhdr}
\usepackage{tikz}

\setlist{nosep, leftmargin=1.2em}

% ═══════════════════════════════════════════════════════════════
% FARBEN
% ═══════════════════════════════════════════════════════════════

\definecolor{physik}{HTML}{2c5aa0}
\definecolor{grau}{HTML}{888888}
\definecolor{hellgrau}{HTML}{f5f5f5}
\definecolor{hinweis}{HTML}{b35c00}

\newcommand{\fachfarbe}{physik}

% ═══════════════════════════════════════════════════════════════
% KOPFZEILE
% ═══════════════════════════════════════════════════════════════

\pagestyle{fancy}
\fancyhf{}
\renewcommand{\headrulewidth}{0.4pt}
\renewcommand{\footrulewidth}{0pt}

\lhead{\small Physik 10E – Waagerechter Wurf}
\rhead{\small\textcolor{physik}{\textbf{Lernbegleiter}}}

% ═══════════════════════════════════════════════════════════════
% BEFEHLE
% ═══════════════════════════════════════════════════════════════

\newcommand{\luecke}[1]{\rule[-2pt]{#1}{0.5pt}}
\newcommand{\ph}[1]{\textcolor{\fachfarbe}{\textbf{#1}}}

\newtcolorbox{formelbox}{
    colback=white,
    colframe=\fachfarbe,
    boxrule=1.2pt,
    arc=0pt,
    left=5pt, right=5pt, top=5pt, bottom=5pt
}

\newtcolorbox{merkbox}[1][]{
    colback=hellgrau,
    colframe=\fachfarbe,
    boxrule=0.8pt,
    arc=0pt,
    left=5pt, right=5pt, top=4pt, bottom=4pt,
    fonttitle=\bfseries\small,
    title=#1
}

\newtcolorbox{aufgabenbox}{
    colback=white,
    colframe=grau,
    boxrule=0.5pt,
    arc=0pt,
    left=5pt, right=5pt, top=5pt, bottom=5pt
}

\newcommand{\schritt}[1]{%
    \vspace{4pt}
    \noindent\textcolor{\fachfarbe}{\rule{0.3cm}{0.3cm}}\hspace{4pt}\textbf{\textcolor{\fachfarbe}{#1}}
    \vspace{2pt}
}

% ═══════════════════════════════════════════════════════════════
% DOKUMENT
% ═══════════════════════════════════════════════════════════════

\begin{document}

\begin{center}
{\Large\bfseries\textcolor{physik}{Lernbegleiter: Waagerechter Wurf}}\\[4pt]
{\small Name: \luecke{5cm} \hfill Datum: \luecke{3cm}}
\end{center}

\vspace{-2pt}
\hrule
\vspace{6pt}

\begin{multicols}{2}
\small

% ─────────────────────────────────────────────────────────────
% TEIL 1: GRUNDLAGEN
% ─────────────────────────────────────────────────────────────

\schritt{1–3: Überlagerungsprinzip}

\begin{merkbox}[Merke]
Waagerechter Wurf = Überlagerung von:
\begin{itemize}[leftmargin=1em]
    \item \textbf{Horizontal:} gleichförmige Bewegung ($v_x$ = konstant)
    \item \textbf{Vertikal:} freier Fall ($a = g$)
\end{itemize}
\vspace{2pt}
$\Rightarrow$ Die horizontale Geschwindigkeit beeinflusst die \textbf{Fallzeit nicht}!
\end{merkbox}

\vspace{4pt}

\schritt{4–5: Formeln}

\begin{formelbox}
\begin{tabular}{@{}ll@{}}
\textbf{Fallzeit:} & $t = \sqrt{\dfrac{2h}{g}}$ \\[12pt]
\textbf{Wurfweite:} & $s_x = v_0 \cdot t$ \\
\end{tabular}

\vspace{4pt}
{\footnotesize Mit $g = 10\,\text{m/s}^2$ (Kopfrechnen)}
\end{formelbox}

\vspace{6pt}

\schritt{6–7: Übung Grundaufgabe}

\begin{aufgabenbox}
$h = 20\,\text{m}$, \quad $v_0 = 15\,\text{m/s}$

\vspace{6pt}
$t = \sqrt{\dfrac{2 \cdot \luecke{1cm}}{10}} = \sqrt{\luecke{1cm}} = $ \luecke{1.2cm} s

\vspace{6pt}
$s_x = \luecke{1cm} \cdot \luecke{1cm} = $ \luecke{1.5cm} m
\end{aufgabenbox}

\vspace{6pt}

% ─────────────────────────────────────────────────────────────
% TEIL 2: AUFPRALL
% ─────────────────────────────────────────────────────────────

\schritt{8–10: Aufprallwerte}

\begin{formelbox}
\begin{tabular}{@{}ll@{}}
\textbf{Vertikale Endgeschw.:} & $v_y = g \cdot t$ \\[8pt]
\textbf{Resultierende:} & $v_{\text{res}} = \sqrt{v_x^2 + v_y^2}$ \\[8pt]
\textbf{Aufprallwinkel:} & $\tan(\alpha) = \dfrac{v_y}{v_x}$ \\
\end{tabular}
\end{formelbox}

\vspace{4pt}

\begin{center}
\begin{tikzpicture}[scale=0.7]
    % Dreieck
    \draw[thick] (0,0) -- (3,0) node[midway, below] {$v_x$};
    \draw[thick] (3,0) -- (3,-2) node[midway, right] {$v_y$};
    \draw[thick, physik] (0,0) -- (3,-2) node[midway, above left] {$v_{\text{res}}$};
    % Winkel
    \draw (0.8,0) arc (0:-34:0.8);
    \node at (1.2,-0.25) {\small$\alpha$};
\end{tikzpicture}
\end{center}

\vspace{2pt}

{\small\textbf{Tangens-Tabelle:}}

\begin{center}
\begin{tabular}{|c|c|c|c|c|c|}
\hline
$\alpha$ & 30° & 37° & 45° & 53° & 60° \\
\hline
$\tan$ & 0,58 & 0,75 & 1,00 & 1,33 & 1,73 \\
\hline
\end{tabular}
\end{center}

\vspace{2pt}

\begin{aufgabenbox}
\textbf{Aufprallwerte} (mit $t = 2\,\text{s}$, $v_0 = 15\,\text{m/s}$):

\vspace{4pt}
$v_y = 10 \cdot 2 = $ \luecke{1.2cm} m/s

\vspace{4pt}
$v_{\text{res}} = \sqrt{15^2 + 20^2} = \sqrt{\luecke{1cm}} = $ \luecke{1.2cm} m/s

\vspace{4pt}
$\tan(\alpha) = \dfrac{20}{15} = $ \luecke{1cm} $\Rightarrow \alpha = $ \luecke{1cm}°
\end{aufgabenbox}

% ─────────────────────────────────────────────────────────────
% SPALTENUMBRUCH
% ─────────────────────────────────────────────────────────────

\columnbreak

% ─────────────────────────────────────────────────────────────
% TEIL 3: ANWENDUNGEN
% ─────────────────────────────────────────────────────────────

\schritt{11: Anwendung – Hilfspaket}

\begin{aufgabenbox}
$h = 80\,\text{m}$, \quad $v_0 = 30\,\text{m/s}$

\vspace{6pt}
\begin{tabular}{@{}ll@{}}
$t = $ & \luecke{2cm} s \\[8pt]
$s_x = $ & \luecke{2cm} m \\[8pt]
$v_y = $ & \luecke{2cm} m/s \\[8pt]
$v_{\text{res}} = $ & \luecke{2cm} m/s \\[8pt]
$\alpha = $ & \luecke{2cm}° \\
\end{tabular}
\end{aufgabenbox}

\vspace{6pt}

\schritt{12: Umkehraufgabe}

\begin{aufgabenbox}
Gegeben: $t = 4\,\text{s}$, \quad $s_x = 100\,\text{m}$

\vspace{6pt}
$v_0 = \dfrac{s_x}{t} = $ \luecke{2cm} m/s

\vspace{6pt}
$h = \dfrac{1}{2} g t^2 = $ \luecke{2cm} m
\end{aufgabenbox}

\vspace{8pt}

% ─────────────────────────────────────────────────────────────
% TEIL 4: ERWEITERUNG
% ─────────────────────────────────────────────────────────────

\schritt{13: Schräger Wurf}

\begin{merkbox}
\textbf{Überlagerung:}
\begin{enumerate}[leftmargin=1.5em]
    \item Gleichförmige Horizontalbewegung
    \item Senkrechter Wurf nach \luecke{1.5cm}
\end{enumerate}

\vspace{4pt}
Bei symmetrischem Wurf: Steigzeit = \luecke{2cm}

\vspace{4pt}
Maximale Wurfweite bei Winkel: \luecke{1.5cm}°

\vspace{4pt}
Bahnform: \luecke{2.5cm}
\end{merkbox}

\vspace{6pt}

\schritt{14: Orbitalmechanik}

\begin{merkbox}
\textbf{Erkenntnis:} Ein Orbit ist ein kontinuierlicher \luecke{3cm}!

\vspace{6pt}
\textbf{Kosmische Geschwindigkeiten:}
\begin{itemize}[leftmargin=1em]
    \item $v_1 \approx$ \luecke{1.2cm} km/s \small(Kreisbahn)
    \item $v_2 \approx$ \luecke{1.2cm} km/s \small(Fluchtgeschw.)
\end{itemize}

\vspace{6pt}
\textbf{Warum stürzt die ISS nicht ab?}

\vspace{2pt}
\luecke{6cm}

\luecke{6cm}
\end{merkbox}

\vspace{8pt}

% ─────────────────────────────────────────────────────────────
% SELBSTCHECK
% ─────────────────────────────────────────────────────────────

\hrule
\vspace{4pt}

{\small\textbf{Selbstcheck:}}

\vspace{2pt}
{\footnotesize
$\square$ Überlagerungsprinzip verstanden \quad
$\square$ Formeln übertragen \\
$\square$ Grundaufgabe gerechnet \quad
$\square$ Aufprallwerte berechnet \\
$\square$ Anwendung + Umkehr gelöst \quad
$\square$ Schräger Wurf \quad
$\square$ Orbit
}

\end{multicols}

\end{document}
